\chapter{Introduction}\label{introduction}

\section{Motivation}\label{motivation}

Large language models (LLMs) are used in software development for more than code completion. Coding agents can inspect repositories, edit multiple files, execute tools, and iteratively respond to test results. \textcite{kwa2025}, for example, propose a task-completion time horizon that relates an agent's success to the time required by a human expert and report an approximately exponential historical increase on their studied software and reasoning tasks. Benchmarks such as SWE-bench have likewise moved evaluation from isolated code completion toward the resolution of real issues drawn from software repositories \autocite{jimenez2024}. These developments expand the scope of implementation work that can be delegated to models and make increasingly autonomous implementation workflows technically plausible.

Evidence for end-to-end productivity is less clear than benchmark progress. In
a randomized controlled trial, experienced open-source developers worked on
repositories they already knew. The early-2025 tools tested by \textcite{becker2025} increased their task completion time. Benchmark capability and observed
delivery speed are thus separate empirical outcomes. Generated code still has
to be placed in repository context, checked against the intended behavior, and
reviewed.

This thesis starts from a narrower observation. When producing a candidate
implementation becomes cheaper, requirements clarification and verification
account for a larger share of the remaining effort. An agent can turn
underspecified text into code before the tacit clarifications of a development
team have occurred. A higher volume of generated changes also increases the
amount of behavior that must be checked. Faster implementation consequently
raises the value of precise requirements and efficient verification, even
though its effect on total engineering time varies between settings.

UI requirements expose this shift clearly because the contract and its evidence
use different modalities. The contract is usually written in natural language;
the evidence appears in visible states and transitions. For example, a
requirement may concern unauthenticated menu access, a filter that remains
selected, or fees displayed in a cart summary. No single source-code location
necessarily establishes any of these behaviors.

Manual review can interpret such behavior but takes time and may vary between
reviewers. Conventional GUI automation provides repeatable interactions when
an executable system, robust selectors, and explicit oracles are available.
\textcite{nass2021} describe practical challenges arising from
the complexity and variability of graphical interfaces. Recorded screenshot
sequences preserve concrete states and their order without requiring a live
deployment. Their evidential scope is narrower: they contain one observed
trajectory and omit other reachable states and hidden implementation
properties.

\section{Problem Statement}\label{problem-statement}

This thesis studies requirements-based review of UI flows. An ordered screenshot sequence records the successive interface states available to the reviewer and is paired with a collection of textual requirements. The screenshots are therefore the evidence representation for the review, rather than the object of the requirements themselves. For each requirement, the system should produce a verification label and point to the screenshot steps that justify its decision. The label set is \texttt{FULFILLED}, \texttt{PARTIALLY\_\allowbreak FULFILLED}, \texttt{NOT\_FULFILLED}, and \texttt{ABSTAIN}. The output may additionally contain smaller requirement claims, claim-level evidence statuses, uncertainty reasons, a natural-language rationale, and bounding boxes that localize visible supporting or contradicting evidence.

The technical problem is to construct a trace from the requirement to the
available evidence. Related work connects user stories with GUI prototypes,
performs agent-driven requirement verification in interactive GUIs, and
combines textual and visual evidence for mobile-app bug-fix verification
\autocite{kretzer2025,kolthoff2025guispector,massenon2026}.
The setting studied here fixes the evidence artifact in advance as an ordered
set of UI states. Verification therefore requires a decision about
observability, identification of the relevant states, localization of evidence
within those states, and a boundary between visible support and inference.

The label boundary is clearest when a requirement explicitly asks for an
observable outcome. In one reviewed benchmark item, the flow shows a configured
eGift card and an \texttt{Add\ to\ Bag} control, but no updated bag or checkout state.
This supports the availability of the submission mechanism, not the requirement
that the configured card be added to the ongoing bag and proceed to checkout.
A \texttt{FULFILLED} label would therefore claim an outcome absent from the recorded
flow.

The label policy follows this evidential boundary. \texttt{FULFILLED} requires support
for all central observable parts. \texttt{NOT\_FULFILLED} requires visible
counter-evidence; an omitted state is insufficient. Support for one material
part combined with a missing, hidden, or ambiguous part can justify
\texttt{PARTIALLY\_\allowbreak FULFILLED}. When the flow supports no reliable positive or negative
decision, the output is \texttt{ABSTAIN}. This last case follows the reject-option
principle for decisions under insufficient information \autocite{hendrickx2024,wen2025}.

The task requires a boundary between visible UI behavior and hidden system
truth. Backend correctness, security guarantees, payment processing, email
delivery, database persistence, global result-set or data completeness,
long-term availability,
and external effects lie outside screenshot evidence unless the requirement
targets their visible representation. A success message can serve as a visible
proxy for the message shown to the user, without establishing the underlying
operation. Predictions and reference labels follow this boundary.

\section{Research Gap and Approach}\label{research-gap-and-approach}

Existing work covers several neighboring tasks. \textcite{kretzer2025} connect user stories with GUI prototypes and study whether a story is represented in a prototype. More directly, GUISpector uses a multimodal agent to operationalize natural-language requirements, explore interactive GUI applications, record verification trajectories, and classify requirements as met, partially met, or unmet \autocite{kolthoff2025guispector}. UI-agent datasets such as Mind2Web represent ordered interaction trajectories for language-conditioned action selection \autocite{deng2023}, while SeeClick studies the localization of interface elements from language instructions \autocite{cheng2024}. Cross-modal software-verification research also combines textual and visual artifacts to assess mobile-app bug fixes \autocite{massenon2026}. Requirements traceability explains why links between statements and artifacts matter, while abstention research explains why a model should sometimes decline a definite decision.

The evidence regime and output contract distinguish the setting studied here.
GUISpector creates evidence by operating an application; this thesis begins with
a previously recorded flow and asks what that fixed, incomplete record
establishes. Its label contract adds \texttt{ABSTAIN} and reserves \texttt{NOT\_FULFILLED} for
visible contradiction, thereby separating insufficient evidence from
requirement violation. The task also differs from action selection and element
localization. Screenshot order matters whenever an action appears in one state
and its result in a later state.

The implementation converts each screenshot into a lightweight representation
and assesses the requirement's UI evaluability. Depending on the experimental
condition, it keeps the requirement intact or decomposes it into claims. A
retriever selects candidate steps, the multimodal verifier assigns claim
statuses, and deterministic rules aggregate them into a requirement label.
Stored intermediate outputs expose retrieval and reasoning errors separately.

``Evidence-first'' names this architecture; its performance remains an empirical
question. In the current experiments, the staged configurations do not
consistently outperform direct whole-flow verification. Their demonstrated
benefit is diagnostic: evidence quality can be measured, and omitted states can
be distinguished from errors made after the correct state was supplied.

\section{Research Questions}\label{research-questions}

The thesis is organized around three research questions:

\textbf{RQ1: How well can multimodal models determine the fulfillment status of written requirements from ordered sequences of UI screenshots?}

RQ1 measures end-to-end agreement with the evidence-bounded four-label
operationalization defined in Section 3.4. The analysis uses accuracy,
macro-F1, per-class measures, confusion matrices, inter-model agreement, and
false fulfillment.

\textbf{RQ2: How do claim decomposition and screenshot selection affect label accuracy, evidence traceability, and cost relative to direct whole-flow verification?}

RQ2 varies claim decomposition and screenshot selection independently. A
controlled matrix compares raw and decomposed requirements under complete-flow
and shared top-k input. Label correctness, evidence overlap, runtime, token usage, and
cost are reported separately.

\textbf{RQ3: Which requirement and evidence patterns cause verification errors, abstentions, or unsafe \texttt{FULFILLED} predictions?}

RQ3 applies a predefined error taxonomy to schema violations, reasoning errors,
retrieval failures, insufficient evidence, and label-boundary disagreements.
The analysis covers universal and comparative wording, late result states,
hidden outcomes, and cross-step persistence. It characterizes observed
abstentions without assigning them a causal safety effect.

\section{Contributions and Scope}\label{contributions-and-scope}

The intended contributions of the thesis are:

\begin{enumerate}
\def\labelenumi{\arabic{enumi}.}
\tightlist
\item
  A problem formulation for verifying textual UI requirements against ordered
  screenshot flows, including a separation between UI evaluability and
  fulfillment.
\item
  An evidence-first prototype that decomposes requirements, retrieves
  screenshot evidence, performs screenshot-grounded claim verification, and
  applies conservative aggregation.
\item
  A reviewed benchmark over 13 Mind2Web-derived flows containing requirement
  labels, claims, evidence steps, rationales, uncertainty reasons, and
  contrastive items.
\item
  An evaluation framework for label quality, false fulfillment, abstention,
  coverage, evidence retrieval, claim matching, runtime, and cost.
\item
  An empirical analysis of the requirement structures and flow patterns behind
  unsafe or uncertain decisions.
\end{enumerate}

Screenshot-step evidence and region-level evidence are evaluated as two
different outputs. Screenshot-step retrieval is evaluated automatically against
the reviewed evidence steps. Region-level grounding is not treated as a cause of
better label accuracy; it is an additional traceability output whose relevance,
sufficiency, and coverage require human evaluation. The repository contains
candidate-mark grounding runs over all 13 flows and the review infrastructure
needed for this evaluation. PURE provides a separate exploratory
document-to-UI consistency analysis; the primary benchmark remains the 13
recorded execution flows.

The empirical scope is limited to 13 flows and does not establish industrial
readiness or broad generalization. The study makes the task and its failure
modes measurable, compares concrete verifier designs, and identifies the claims
supported by screenshot evidence.

The remainder of the thesis first establishes the conceptual and technical
foundations, formalizes the verification task, describes the implemented
pipeline, and then presents the benchmark, evaluation, results, limitations, and
conclusions.
